% =====================================================================
\chapter{Giới thiệu}\label{chap:intro}
% =====================================================================

\section{Mục đích tài liệu}
Tài liệu này là \textbf{tài liệu thiết kế hệ thống} (System Design Document) cho nền tảng dữ liệu
\emph{Marketplace Price Intelligence}. Nó trả lời năm nhóm câu hỏi mà một kỹ sư dữ liệu mới tham gia,
một người vận hành hoặc một người duyệt kỹ thuật cần biết:

\begin{enumerate}
  \item Hệ thống giải quyết bài toán gì, phạm vi tới đâu, yêu cầu nào là bắt buộc (Chương~\ref{chap:req}).
  \item Kiến trúc gồm những tầng nào, chạy ở đâu, vì sao chọn như vậy (Chương~\ref{chap:arch}, Phụ lục~\ref{app:adr}).
  \item Dữ liệu được thu thập, lưu trữ, mô hình hoá và biến đổi thế nào qua từng tầng (Chương~\ref{chap:ingest}--\ref{chap:processing}).
  \item Có những pipeline nào, lịch chạy, đầu vào/ra, cách chạy lại khi lỗi; cổng chất lượng chặn gì (Chương~\ref{chap:pipeline}--\ref{chap:serving}).
  \item Hệ thống được triển khai, giám sát, sao lưu và khôi phục ra sao; đã đo được gì; còn rủi ro gì (Chương~\ref{chap:deploy}--\ref{chap:risk}).
\end{enumerate}

\section{Phạm vi tài liệu}
Tài liệu mô tả trạng thái mã nguồn tại \DocBaseline, tức nhánh \code{feat/discount-credibility}
đã merge thêm \code{phase-11/databricks}. Phạm vi gồm toàn bộ package \code{ecomdata}, cấu hình Docker Compose,
DAG Airflow, job Databricks, DDL PostgreSQL và cấu hình giám sát.

Ngoài phạm vi: mã thử nghiệm cũ trên dữ liệu Kaggle (thư mục \code{batch_layer/}, \code{speed_layer/},
\code{serving_layer/} ở gốc repository, giữ lại làm di sản), kế hoạch các phase đã lưu trữ trong \code{docs/archive/}.

\section{Đối tượng đọc}
\begin{xltabular}{\linewidth}{@{}L{40mm}Y@{}}
  \toprule
  \textbf{Đối tượng} & \textbf{Nên đọc} \\
  \midrule
  Người hướng dẫn, hội đồng & Chương~\ref{chap:req}, \ref{chap:arch}, \ref{chap:perf}, \ref{chap:risk} \\
  Kỹ sư dữ liệu phát triển & toàn bộ, đặc biệt Chương~\ref{chap:storage}--\ref{chap:quality} \\
  Người vận hành (on-call) & Chương~\ref{chap:pipeline}, \ref{chap:deploy}, \ref{chap:ops}; \code{docs/RUNBOOK.md} \\
  Người dùng dữ liệu (analyst) & Mục~\ref{sec:marts} (danh mục mart), Chương~\ref{chap:serving} \\
  \bottomrule
\end{xltabular}

\section{Cách đọc tài liệu}
Các chương đi theo đường đi của dữ liệu: từ trang của sàn (thu thập), qua lake (Bronze, Silver, Gold),
qua hai đường xử lý (batch và speed), tới nơi phục vụ (PostgreSQL, Elasticsearch, Redis) và người dùng
(Superset, Kibana). Mỗi chương ghi rõ \emph{nguồn chân lý} là file nào trong mã nguồn.
Thuật ngữ riêng của hệ thống được định nghĩa ở Mục~\ref{sec:glossary}; nên đọc trước khi vào Chương~\ref{chap:storage}.

\section{Thuật ngữ}\label{sec:glossary}
\begin{xltabular}{\linewidth}{@{}L{38mm}Y@{}}
  \caption{Thuật ngữ dùng trong hệ thống}\label{tab:glossary}\\
  \toprule
  \textbf{Thuật ngữ} & \textbf{Nghĩa trong hệ thống này} \\
  \midrule
  \endfirsthead
  \toprule \textbf{Thuật ngữ} & \textbf{Nghĩa} \\ \midrule
  \endhead
  offer & một listing của một sàn: (\code{source_code}, \code{platform_listing_id}) \ra\ \code{offer_id} \\
  observation (sighting) & một lần crawler thấy một offer trên một trang, ở một thời điểm \\
  record & một dòng nguồn nguyên vẹn (envelope \code{marketplace-record.v1}), chưa chuẩn hoá \\
  raw artifact & toàn bộ phản hồi của một lần tải trang (HTML và các XHR), lưu nguyên byte trong Bronze/raw \\
  run & một lượt batch, khoá bởi \code{run_id} (\code{mp-YYYYMMDDTHHMMZ}) và \code{as_of} \\
  \code{as_of} & mốc thời gian dữ liệu mà run được phép thấy; mart và quality cắt tại đây \\
  cửa sổ ngày & tập partition \code{observed_date} mà một run được dựng lại ở Silver (ngày của \code{as_of} và \code{late_days} ngày trước) \\
  star / \code{dw} & 7 dimension, 1 bridge và 6 fact (Delta) dùng chung cho mọi run \\
  mart & bảng phục vụ theo run (12 bảng); bản được promote được sao vào PostgreSQL \code{cache.*} \\
  manifest & file JSON mô tả một run: dataset, kết quả quality, version Delta của từng bảng star \\
  pointer & \code{current.json}: bản sao byte của manifest của run đang được phục vụ \\
  promote & thay pointer bằng compare-and-swap để run mới trở thành run được phục vụ \\
  counter & số mà sàn hiển thị (đã bán, số đánh giá). Thay đổi của counter \textbf{không phải} đơn hàng \\
  frontier & hàng đợi việc crawl trong \code{audit.crawl_frontier} \\
  tier & mức độ ưu tiên/chu kỳ crawl của một target: \code{ACTIVE}, \code{HOT}, \code{NORMAL}, \code{COLD}, \code{SWEEP} \\
  quarantine & bảng Silver giữ các dòng không hợp lệ kèm lý do, thay vì bỏ \\
  DLQ & dead-letter queue: topic Kafka nhận các bản tin không xử lý được của speed \\
  organic & dòng sản phẩm không phải quảng cáo; mart và speed chỉ xét dòng organic \\
  \bottomrule
\end{xltabular}

\section{Tài liệu liên quan}
\begin{xltabular}{\linewidth}{@{}L{62mm}Y@{}}
  \toprule
  \textbf{Tài liệu} & \textbf{Nội dung} \\
  \midrule
  \code{docs/system/README.md}, \code{01}--\code{05} & bản Markdown của tài liệu này, kèm spec đầy đủ (bảng node, cạnh) của từng sơ đồ \\
  \code{docs/adr/0001}--\code{0009} & các quyết định kiến trúc \\
  \code{docs/model/star-schema.dbml} & cột đầy đủ của star schema (sinh ra \code{.mmd}, \code{.puml}, \code{.sql}, PNG) \\
  \code{docs/RUNBOOK.md} & quy trình vận hành và xử lý sự cố \\
  \code{docs/OBSERVABILITY.md} & giám sát: metrics, log, dashboard, cảnh báo \\
  \code{docs/DATABRICKS_READINESS.md} & mức sẵn sàng của phần chạy trên Databricks \\
  \code{docs/DANH_GIA_HE_THONG_20261113.md} & đánh giá hệ thống ngày 13/11/2026 \\
  \code{docs/evidence/fullrun-20261114/KET-QUA.md} & kết quả lượt chạy toàn luồng 13--14/11/2026 \\
  \bottomrule
\end{xltabular}
